\chapter*{Abstract}
\addcontentsline{toc}{chapter}{Abstract}

Large-scale manufacturers such as IKEA of Sweden rely on hundreds of standardized communication
labels—product identity cards, packaging inserts, logistics tags, and regulatory compliance
markers—that must comply with strict layout and content specifications before every print run.
Manual quality assurance of these templates is slow, inconsistent, and expensive at production
volume. This thesis presents the design and evaluation of an automated label validation
pipeline that combines multi-engine Optical Character Recognition (\ac{OCR}), rule-based
field validation, GS1 barcode decoding, and a scoring-based label-type classifier to determine
whether a given label image conforms to its target specification.

The system was built in close collaboration with IKEA of Sweden, whose internal template
catalogue of 34 distinct label types formed the basis of the ground-truth rule schemas.
A dataset of 238 labelled images—131 compliant, 107 non-compliant—was assembled by
collecting real production-stage label PDFs and programmatically introducing controlled
violations. Each image was annotated at the field level, recording which required elements
were present, their extracted text values, barcode payloads, and the compliance outcome.

Three \ac{OCR} engines were benchmarked: Tesseract~5.3, EasyOCR~1.7, and a late-fusion
ensemble of the two. The ensemble achieved a character-level accuracy of~89.1\% and a
word-level accuracy of~83.2\%, outperforming either engine alone. Barcode decoding
reached 94.4\% for ITF-14 symbols and 88.7\% for Data\-Matrix symbols using the
pyzbar and pylibdmtx libraries respectively. Label-type identification—mapping an
unknown image to one of 34 schema candidates—attained a top-1 accuracy of~86.1\%
and a top-3 accuracy of~94.4\% on the held-out test set. Field extraction achieved
a macro-averaged F1 score of~0.866 across ten key fields, with the highest scores
on structured identifiers such as the IKEA supplier name (F1~=~0.957) and article
number (F1~=~0.934), and lower scores on free-text fields such as physical dimensions
(F1~=~0.804). The rule-based validator that applies schema-specific checks to the
extracted fields reached a precision of~88.3\% and a recall of~84.7\% in flagging
non-compliant labels, yielding an F1 of~0.865.

The results demonstrate that a hybrid pipeline combining deterministic OCR, GS1
barcode parsing, and rule-based validation can automate the majority of IKEA's label
quality-assurance workload with high reliability, reducing average manual review time
from an estimated~8 minutes per label to under~2 seconds of automated processing.
Remaining error modes are concentrated in low-resolution scans, rotated text zones,
and partially occluded barcodes, pointing to clear directions for future work.

\vspace{1em}
\noindent\textbf{Keywords:}
optical character recognition, rule-based validation, document image analysis,
label compliance, GS1 barcodes, template classification, field extraction,
IKEA, computer vision, explainability.
\clearpage
